% BHU PYQ -- Transcribed & Corrected
% Paper: MATMJ-32 / MATMJ32: Analysis (Mid Term)
% Exam: B.Sc. (Hons.) Semester III Mid-Term Examination, 2025-26 (NEP)
% Department: Mathematics

\documentclass[11pt,a4paper]{article}
\usepackage[a4paper,top=2.5cm,bottom=2.5cm,left=2.8cm,right=2.8cm,headheight=14pt]{geometry}
\usepackage{amsmath,amssymb}
\usepackage[T1]{fontenc}
\usepackage[utf8]{inputenc}
\usepackage{lmodern,microtype,fancyhdr,enumitem,hyperref,lastpage,parskip}

\hypersetup{
    colorlinks=true,
    urlcolor=black,
    linkcolor=black,
    pdftitle={MATMJ32: Analysis Mid-Term -- BHU B.Sc. Sem III 2025-26 (NEP)}
}

\pagestyle{fancy}
\fancyhf{}
\renewcommand{\headrulewidth}{0.4pt}
\renewcommand{\footrulewidth}{0.4pt}
\fancyhead[L]{\small\textit{Banaras Hindu University}}
\fancyhead[R]{\small\textit{MATMJ-32: Analysis (Mid Term)}}
\fancyfoot[C]{\small Page \thepage\ of \pageref{LastPage}}

\newlist{parts}{enumerate}{2}
\setlist[parts,1]{label=(\roman*),leftmargin=*,itemsep=4pt,topsep=2pt}
\newcommand{\pts}[1]{\hfill\textit{[#1]}}

\begin{document}

\begin{center}
  {\large\bfseries BANARAS HINDU UNIVERSITY}\\[2pt]
  {\normalsize B.Sc. (Hons.) --- Semester III Mid-Term Examination, 2025-26 (NEP)}\\[6pt]
  \rule{0.8\linewidth}{0.5pt}\\[6pt]
  {\large\bfseries MATHEMATICS}\\[4pt]
  {\large\bfseries Paper Code: MATMJ-32 : Analysis (Mid Term)}\\[4pt]
  {\normalsize Time: 1 Hour\hfill Maximum Marks: 20}
\end{center}

\rule{\linewidth}{0.4pt}

\smallskip
\noindent\textit{Instructions: Attempt any \textbf{5} questions. All questions carry equal marks (4 marks each).}

\bigskip

\noindent\textbf{Question 1.} What are Peano's Axioms? State and prove the Law of Trichotomy for $\mathbb{N}$, OR prove that $m \cdot (n + p) = m \cdot n + m \cdot p$ for all $m, n, p \in \mathbb{N}$.\pts{4}

\medskip\rule{\linewidth}{0.2pt}\medskip

\noindent\textbf{Question 2.} What are the common algebraic and order properties shared between the set of real numbers $\mathbb{R}$ and the set of rational numbers $\mathbb{Q}$? Mention the completeness property (Least Upper Bound Axiom) that distinguishes $\mathbb{R}$ from $\mathbb{Q}$. Illustrate this difference with a concrete example.\pts{4}

\medskip\rule{\linewidth}{0.2pt}\medskip

\noindent\textbf{Question 3.} What do you mean by a convergent sequence? Prove that a bounded sequence $(x_n)$ is convergent if and only if $\limsup x_n = \liminf x_n$.\pts{4}

\medskip\rule{\linewidth}{0.2pt}\medskip

\noindent\textbf{Question 4.} Differentiate between an absolutely convergent series and a conditionally convergent series. Provide standard examples for both types. Test the convergence of the series $\sum_{n=1}^\infty \frac{\sin n\theta}{n^2}$.\pts{4}

\medskip\rule{\linewidth}{0.2pt}\medskip

\noindent\textbf{Question 5.} If $f \colon [0, a] \to \mathbb{R}$ is defined by $f(x) = x^2$ with $a > 0$, prove that $f \in \mathcal{R}[0, a]$ (Riemann integrable) and evaluate $\int_0^a f(x) \, dx = \frac{a^3}{3}$.\pts{4}

\medskip\rule{\linewidth}{0.2pt}\medskip

\noindent\textbf{Question 6.} Prove Cauchy's Criterion for Integrability: A bounded function $f \colon [a, b] \to \mathbb{R}$ is Riemann integrable if and only if for any $\epsilon > 0$, there exists a partition $P$ of $[a, b]$ such that $U(P, f) - L(P, f) < \epsilon$.\pts{4}

\vfill
\rule{\linewidth}{0.4pt}
\begin{center}\small
\href{https://bhu-exam.vercel.app/}{Click here to visit our website}\\[4pt]
\href{https://me-aryan.vercel.app/}{Click here to visit our contributor}
\end{center}

\end{document}
